\section{Data and preprocessing}
\label{data}

\subsection{Exploration}

Each ad has five free-text fields (title, company profile, description, requirements, benefits),
categorical metadata (employment type, industry, function, location, required experience and
education) and three binary flags (telecommuting, company logo, screening questions); real URLs,
emails and phone numbers are already replaced by hashed placeholders. Each step of the exploration
(notebook \texttt{01}, on the 17,880 raw ads) answered a question that fixed a later design
decision, as summarised in Table~\ref{tab:eda}.

The 866 fraudulent ads (4.84\%) make accuracy useless, so we use PR-AUC, choose the decision
threshold on validation, and weight the loss rather than resample. 281 rows are exact duplicates
and 4,552 rows share their description with another row (up to 380 ads per description), so the
data must be split by description, not by row. Missing fields are a signal rather than noise
(Figure~\ref{fig:flags}): a company profile is missing in 67.8\% of fraudulent ads against 16.0\% of
real ones, and fraudulent ads more often lack a logo and contain money words or contact details.
Finally, some countries have no fraud at all over hundreds of ads (Greece, 940 ads), which is more
likely a property of how EMSCAD was collected than of the job markets; we keep the feature but
check in Section~\ref{explain} how much the models rely on it.

\begin{figure}[tb]
\centering
\includegraphics[width=0.72\linewidth]{eda_flags.pdf}
\caption{Share of ads in which each flag is set, real versus fraudulent (17,880 raw ads).}
\label{fig:flags}
\end{figure}

\begin{table}[tb]
\centering
\small
\begin{tabular}{>{\raggedright\arraybackslash}p{0.40\linewidth}>{\raggedright\arraybackslash}p{0.52\linewidth}}
\toprule
\textbf{Finding} & \textbf{Decision} \\
\midrule
4.84\% of ads are fraudulent. & PR-AUC; threshold chosen on validation; class weights. \\
281 exact duplicates; 4,552 rows share a description. & Drop duplicates; split by description. \\
Missing fields are far more common in fraud. & Keep a missing-value flag per field. \\
Some countries have no fraud over hundreds of ads. & Keep country, check reliance on it. \\
\bottomrule
\end{tabular}
\caption{Exploration findings and the design decision each one drives.}
\label{tab:eda}
\end{table}

\subsection{Preprocessing and splits}

Dropping the exact duplicates leaves 17,599 ads, 856 of them fraudulent, and each ad is assigned
the hash of its normalised description as its group. \texttt{StratifiedGroupKFold} on these groups
gives a 60/20/20 train/validation/test split (10,559 / 3,520 / 3,520 ads, with 514 / 171 / 171
frauds) in which no description appears in two parts. We repeat this for five split seeds and save
the row indices, so every notebook and machine uses the same rows.

The five text fields are joined into one column, and a \texttt{ColumnTransformer}, fitted on the
training rows only, builds the feature matrix of Table~\ref{tab:features}. TF-IDF weighs each word
by how often it occurs in the ad and how rare it is across ads.

\begin{table}[tb]
\centering
\small
\begin{tabular}{l>{\raggedright\arraybackslash}p{0.55\linewidth}r}
\toprule
\textbf{Block} & \textbf{Transformer} & \textbf{Columns} \\
\midrule
Text & TF-IDF, word 1- and 2-grams & 20,000 \\
Metadata & One-hot categories, ordered experience and education levels & 112 \\
Numeric & Text lengths, exclamation marks, share of capitals (scaled) & 4 \\
Flags & Logo, questions, telecommuting and 9 engineered flags & 12 \\
\bottomrule
\end{tabular}
\caption{Feature matrix (default settings, seed-0 training rows).}
\label{tab:features}
\end{table}
